% Material migrated from the former appendix into the main body for IPDPS.
% All added prose is review-marked so coauthors can inspect it before finalization.

\anurag{Appendix - B Methodology: Data Collection to predict communication energy}

\review{\noindent\textbf{Parallelism-specific communication measurements.}
The communication-energy measurements used to train \sys{} are collected differently for the three parallelism strategies because the communication event boundaries are different. For tensor parallelism, we repeatedly profile the ring AllReduce operation rather than relying on a single trace. This is important because the ranks do not always reach the collective at exactly the same time: one GPU may lead or lag another because of runtime scheduling, cache behavior, or memory-access variation. The measured AllReduce interval therefore contains both link activity and synchronization/waiting energy. By collecting repeated executions, \sys{} captures the distribution of these outcomes and avoids training the communication predictor on one potentially unrepresentative synchronization instance.}

\review{For pipeline parallelism, communication is isolated at every stage boundary. We use \texttt{nsys} timestamps to identify the completion of the last forward kernel on stage $i$ and the start of the first forward kernel on stage $i+1$. The interval between these two events is attributed to the point-to-point activation transfer, and system power is integrated over the same interval to obtain the corresponding communication energy. This measurement procedure directly associates the transferred activation tensor with the energy consumed while moving it between the producing and consuming stages.}

\review{For data parallelism, each replica completes its local forward pass independently and a final AllGather is used to collate the outputs. We therefore profile the tail-end interval beginning when the output layer completes and ending when the AllGather has finished. Because there is only one such output-collation event per inference request in our setup, the communication boundary is more straightforward than in tensor or pipeline parallelism and typically accounts for a smaller fraction of the run. Across all three strategies, bytes sent/received are derived from the communicated tensor size and communication pattern, while effective bandwidth is computed from observed transfer volume and duration. These measurements are collected offline and are aligned with the system power trace used as ground truth.}

\anurag{Appendix - C Methodology: Features}

\review{\noindent\textbf{Feature set.}
Table~\ref{tab:features} contains the complete input feature set used by the module-level predictors. The resource-utilization group captures system and device activity during a module execution, including CPU/GPU utilization, memory utilization, clock frequencies, and memory usage. Execution features capture the workload configuration and observed cost of the run, including batch size, sequence length, FLOPs per token, execution time, NVML-reported GPU energy, and GPU count. Model-structure features such as feed-forward dimension, number of Transformer blocks, hidden size, attention-head count, and key--value-head count expose architectural differences that affect tensor shapes and the amount of work assigned to each rank. For communication modules, \sys{} additionally records transferred bytes and effective bandwidth. This separation is important because the same utilization value can correspond to different energy behavior when tensor shapes, model structure, or communication volume differ.}

\review{The feature design also addresses the scalability concern raised for multi-GPU prediction. Rather than learning a separate set of coefficients for each physical GPU, \sys{} summarizes runtime features across the participating GPUs using mean, standard deviation, minimum, and maximum. The aggregate captures the overall level of activity while the spread and extrema expose imbalance and straggler behavior. These runtime aggregates are then combined with the static structural and communication features before training the module-specific regressors.}

\anurag{Appendix - D Methodology: Details of predicting communication energy}

\review{\noindent\textbf{Communication-volume and effective-bandwidth calibration.}
Section~\ref{ss:design-communication} uses the $\alpha$--$\beta$ relation to predict communication latency. For a communication event $e$, the per-GPU volume is $V_e=B_e^{\mathrm{out}}+B_e^{\mathrm{in}}$. We calibrate the model offline for every communication operation type $o$ and GPU count $p$ by sweeping communication volumes and recording the observed duration $T_e^{\mathrm{obs}}$. Fitting $T_e^{\mathrm{obs}}=\alpha_{o,p}+\beta_{o,p}V_e$ separates a fixed startup/synchronization component, $\alpha_{o,p}$, from the message-size-dependent transfer component, $\beta_{o,p}$. The corresponding effective bandwidth is $B^{\mathrm{eff}}_{o,p}=1/\beta_{o,p}$. This calibration is performed once per hardware/topology and is then reused by the predictor.}

\review{For ring AllReduce, the communication volume follows directly from the two ring phases. A tensor of size $S_e$ is partitioned across $p$ GPUs; each GPU sends and receives $(p-1)S_e/p$ bytes over ReduceScatter and the same amount over AllGather. Consequently, the total per-GPU volume used by Eq.~\ref{eq:comm-latency-main} is $V_e^{AR}=2(p-1)S_e/p$. Point-to-point pipeline transfers use the activation tensor size for the corresponding stage boundary, while data-parallel AllGather volume is derived from the final output tensor and the number of replicas. After predicting $T_e^{\mathrm{comm}}$, \sys{} obtains communication energy from calibrated communication power times predicted communication duration and composes this value with the compute-module estimates in the multi-level regressor.}

\anurag{Appendix - E Methodology: Training}

\review{\noindent\textbf{Training protocol.}
The regressors are trained hierarchically: \sys{} first learns module-level predictors and then composes their outputs to obtain model-level energy. We use 3-fold cross-validation at the module level with a 70\%/30\% train/test split. Each module type (e.g., Self-Attention, MLP, or AllReduce) is sampled $10{,}000$ times to build a robust dataset. A sample is an aggregate measurement formed from repeated executions of that module under a fixed model and GPU configuration; in our profiling campaign each sample aggregates $100{,}000$ module executions. For every sample, the runtime features from all participating GPUs are reduced to their mean, standard deviation, minimum, and maximum values and then combined with the structural features described in Table~\ref{tab:features}.}

\review{To expose the regressors to both low- and high-throughput regimes, experiments span batch sizes 8, 16, 32, and 64 and output sequence lengths 512 and 1024. We use DeepSpeed~\cite{deepspeed} to execute the tensor-, pipeline-, and data-parallel inference runs. At model level, the predicted energies of Attention, MLP, communication, normalization, embedding, and other constituent modules are composed through the multi-level regression in Eq.~\ref{eqn:end2end-predicted-sum}. For the leave-one-variant-out experiments, one complete model variant is excluded from training; module-level predictors learned from the remaining variants in the same family are then composed to predict the held-out model. This protocol directly tests whether the model-tree decomposition and structural features transfer across model scales rather than merely interpolating among samples from the same variant.}
